\noindent This supplementary material records the exact control-panel proofs,
the estimator definitions that bridge theory to observables, and the gap-aware
diagnostic tables retained in the approved manuscript architecture. It
deliberately omits comparator-centered appendices and theorem-forward material
from earlier drafts because those objects are not load-bearing in the final
design.

\section*{S1. Exact control-panel propositions}
\SupplementTheoryOnlyNote

The main text relies on only two exact propositions. They are restated and
proved here.

\begin{proposition}[Common-risk-set factorization]
Fix a horizon \(h>0\). Define
\[
\begin{aligned}
R_h(m)
&:=\mathbf{1}\!\left\{\substack{C_m\ge h\\ \text{or } T_m\le h}\right\},\\
\pi_h
&:=\mathbb{P}(\delta_m=1\mid R_h(m)=1),\\
\phi_h
&:=\mathbb{P}\!\left(T_m\le h \,\middle|\, \substack{\delta_m=1\\ R_h(m)=1}\right),\\
q_h
&:=\mathbb{P}(T_m\le h \mid R_h(m)=1).
\end{aligned}
\]
If \(\mathbb{P}(R_h(m)=1)>0\) and
\(\mathbb{P}(\delta_m=1\mid R_h(m)=1)>0\), then
\[
q_h=\pi_h\phi_h.
\]
\end{proposition}

\begin{proof}
Conditional on \(R_h(m)=1\), the event \(\{T_m\le h\}\) implies
\(\{\delta_m=1\}\): if the first direct reply arrives by \(h\), then a direct
reply is observed before censoring. Therefore,
\begin{align*}
\mathbb{P}(T_m\le h\mid R_h(m)=1)
&=
\mathbb{P}(T_m\le h,\ \delta_m=1 \\
&\hspace{2.5em}\mid R_h(m)=1).
\end{align*}
Applying the chain rule on the right-hand side gives
\begin{gather*}
\mathbb{P}(T_m\le h,\ \delta_m=1\mid R_h(m)=1)\\
=\mathbb{P}(\delta_m=1\mid R_h(m)=1)\times \\
\mathbb{P}\!\left(T_m\le h \,\middle|\, \substack{\delta_m=1\\ R_h(m)=1}\right).
\end{gather*}
which is exactly \(q_h=\pi_h\phi_h\).
\end{proof}

\begin{proposition}[One-margin gains and cost-normalized priority]
Let \(q_h=\pi_h\phi_h\). If only the incidence margin changes by
\(\Delta \pi_h\) while \(\phi_h\) is held fixed, then
\[
\Delta q_h^{(\pi)}=\phi_h\,\Delta \pi_h.
\]
If only the conditional-timing margin changes by \(\Delta \phi_h\) while
\(\pi_h\) is held fixed, then
\[
\Delta q_h^{(\phi)}=\pi_h\,\Delta \phi_h.
\]
For positive intervention costs \(c_\pi\) and \(c_\phi\), the corresponding
cost-normalized local gains are
\[
I_\pi:=\frac{\phi_h\,\Delta \pi_h}{c_\pi},
\qquad
I_\phi:=\frac{\pi_h\,\Delta \phi_h}{c_\phi},
\]
so local priority favors incidence lift if and only if \(I_\pi>I_\phi\).
\end{proposition}

\begin{proof}
If only incidence changes, then
\[
q_h+\Delta q_h^{(\pi)}
=
(\pi_h+\Delta \pi_h)\phi_h
=
\pi_h\phi_h+\phi_h\Delta \pi_h,
\]
hence \(\Delta q_h^{(\pi)}=\phi_h\Delta \pi_h\). If only conditional timing
changes, then
\[
q_h+\Delta q_h^{(\phi)}
=
\pi_h(\phi_h+\Delta \phi_h)
=
\pi_h\phi_h+\pi_h\Delta \phi_h,
\]
hence \(\Delta q_h^{(\phi)}=\pi_h\Delta \phi_h\). Dividing by positive costs
\(c_\pi\) and \(c_\phi\) yields the indices \(I_\pi\) and \(I_\phi\), so the
preferred local margin is incidence if and only if \(I_\pi>I_\phi\).
\end{proof}

\section*{S2. Observable bridge and primary estimators}
\SupplementMixedEvidenceNote

Let \(\mathcal{M}\) denote the realized cohort of non-root parent comments.
For \(m\in\mathcal{M}\) and horizon \(h\), define
\[
\begin{aligned}
Y_h(m)&:=\mathbf{1}\{\delta_m=1,\ T_m\le h\},\\
R_h(m)&:=\mathbf{1}\!\left\{\substack{C_m\ge h\\ \text{or } T_m\le h}\right\}.
\end{aligned}
\]
The sample control-panel estimators are
\begin{align*}
\widehat q_h
&=
\frac{\sum_{m\in\mathcal{M}} Y_h(m)}
     {\sum_{m\in\mathcal{M}} R_h(m)},
\\
\widehat \pi_h
&=
\frac{\sum_{m\in\mathcal{M}} \mathbf{1}\{\delta_m=1\}R_h(m)}
     {\sum_{m\in\mathcal{M}} R_h(m)},
\\
\widehat \phi_h
&=
\frac{\sum_{m\in\mathcal{M}} Y_h(m)}
     {\sum_{m\in\mathcal{M}} \mathbf{1}\{\delta_m=1\}R_h(m)}.
\end{align*}
By construction,
\[
\widehat q_h=\widehat \pi_h\,\widehat \phi_h.
\]
Earlier drafts denoted \(\widehat q_h\) by \(\widehat p_h\). In the approved
architecture, \(\widehat q_h\) is the preferred notation and
\(\widehat p_{\mathrm{obs}}\) is reserved for the descriptive ever-reply share
\[
\widehat p_{\mathrm{obs}}
=
\frac{1}{|\mathcal{M}|}\sum_{m\in\mathcal{M}} \mathbf{1}\{\delta_m=1\}.
\]

All numerical readouts in this appendix come from the canonical
SimulaMet-backed evidence. A larger-archive Moltbook replication remains part
of the study design, but it is not yet a reported empirical result in the
current manuscript.

\subsection*{S2.1 Follow-up-standardized incidence by group}

Table \ref{tab:supp-horizon-incidence-groups} summarizes the primary
follow-up-standardized readouts at \(5\) minutes and \(1\) hour. Claimed versus
unclaimed is retained as descriptive heterogeneity only.

\begin{table}[h]
\centering
\caption{Follow-up-standardized incidence by group. The primary control-panel
throughput quantities are \(q_{5\mathrm{m}}\) and \(q_{1\mathrm{h}}\); earlier
drafts denoted the same quantities by \(p_h\). The in-window ever-reply share
\(p_{\mathrm{obs}}\) is descriptive only. Claimed/unclaimed excludes parents
with missing author identifier (\(n=906\)).}
\label{tab:supp-horizon-incidence-groups}
\small
\begin{tabular}{@{}llrrrr@{}}
\toprule
\textbf{Family} & \textbf{Group} & \textbf{Parents} & \textbf{\(q_{5\mathrm{m}}\) \%} & \textbf{\(q_{1\mathrm{h}}\) \%} & \textbf{\(p_{\mathrm{obs}}\) \%} \\
\midrule
overall & Overall & 223,316 & 9.42 & 9.82 & 9.60 \\
claimed\_status & Claimed & 20,667 & 18.95 & 19.56 & 19.23 \\
claimed\_status & Unclaimed & 201,743 & 8.48 & 8.86 & 8.65 \\
submolt\_category & Social/Casual & 189,765 & 10.80 & 11.21 & 10.95 \\
submolt\_category & Other & 16,396 & 1.98 & 2.29 & 2.26 \\
submolt\_category & Builder/Technical & 8,396 & 1.42 & 1.74 & 1.72 \\
submolt\_category & Philosophy/Meta & 5,831 & 1.32 & 1.70 & 1.80 \\
submolt\_category & Spam/Low-Signal & 2,495 & 0.88 & 0.93 & 0.88 \\
submolt\_category & Creative & 433 & 1.62 & 1.71 & 1.62 \\
\bottomrule
\end{tabular}
\end{table}

\section*{S3. Coverage-gap discipline}
\SupplementMixedEvidenceNote

The canonical comment stream contains one 41.7-hour interruption from
2026-01-31 10:37:53Z to 2026-02-02 04:20:50Z. The approved design therefore
treats gap handling as part of the identification strategy rather than as a
late robustness add-on.

\subsection*{S3.1 Gap-disambiguation evidence}

Table \ref{tab:supp-gap-disambiguation} shows that the interruption is
comment-stream-specific in the archive: comments disappear over the interval,
while posts, snapshots, and word-frequency records continue to arrive.

\begin{table}[h]
\centering
\caption{Comment-gap disambiguation evidence across raw archive tables. The
comment interval gap is 2026-01-31 10:37:53Z to 2026-02-02 04:20:50Z.}
\label{tab:supp-gap-disambiguation}
\small
\begin{tabular}{@{}lrrr@{}}
\toprule
\textbf{Table} & \textbf{Records} & \textbf{Records in comment-gap interval} & \textbf{Max inter-event gap (h)} \\
\midrule
comments & 226,173 & 0 & 41.72 \\
posts & 119,677 & 38,166 & 11.20 \\
snapshots & 114 & 39 & 2.28 \\
word\_frequency & 15,346 & 5,039 & 3.00 \\
\bottomrule
\end{tabular}
\end{table}

\subsection*{S3.2 Window robustness}

Table \ref{tab:supp-gap-window-robustness} reports the two-part decomposition
across the full window, contiguous pre-gap and post-gap windows, and the two
prespecified gap-overlap exclusions.

\begin{table}[h]
\centering
\caption{Two-part decomposition robustness across contiguous windows and
gap-overlap exclusions.}
\label{tab:supp-gap-window-robustness}
\scriptsize
\begin{tabular}{@{}lrrrrrrr@{}}
\toprule
\textbf{Scenario} & \textbf{Parents} & \textbf{Replies} & \textbf{Incidence \%} & \textbf{\(t_{50}\) (s)} & \textbf{\(t_{90}\) (s)} & \textbf{\(\Pr(\mathrm{reply}\le 30\mathrm{s})\) \%} & \textbf{\(\Pr(\mathrm{reply}\le 5\mathrm{m})\) \%} \\
\midrule
Full window & 223,316 & 21,430 & 9.60 & 4.55 & 50.05 & 8.47 & 9.41 \\
Pre-gap contiguous & 2,856 & 30 & 1.05 & 513.48 & 1699.77 & 0.04 & 0.35 \\
Post-gap contiguous & 220,460 & 21,400 & 9.71 & 4.55 & 48.66 & 8.58 & 9.53 \\
Exclude gap overlap (6h) & 220,460 & 21,400 & 9.71 & 4.55 & 48.66 & 8.58 & 9.53 \\
Exclude gap overlap (24h) & 220,460 & 21,400 & 9.71 & 4.55 & 48.66 & 8.58 & 9.53 \\
\bottomrule
\end{tabular}
\end{table}

\subsection*{S3.3 Horizon-standardized risk sets}

Table \ref{tab:supp-gap-horizon-standardized} reports the explicit risk-set
denominators used for the horizon-standardized readouts. Short-horizon values
remain stable, while the one-hour rate rises modestly under standardization,
consistent with right-censoring near the observation boundary rather than with a
slow conditional-timing regime.

\begin{table}[h]
\centering
\caption{Horizon-standardized reply probabilities using explicit risk sets.}
\label{tab:supp-gap-horizon-standardized}
\small
\begin{tabular}{@{}llrr@{}}
\toprule
\textbf{Scenario} & \textbf{Horizon} & \textbf{Risk-set \(n\)} & \textbf{\(\Pr(\mathrm{reply}\le t)\) \%} \\
\midrule
Full window & 30s & 223,312 & 8.47 \\
Full window & 5m & 223,102 & 9.42 \\
Full window & 1h & 217,472 & 9.82 \\
Pre-gap contiguous & 30s & 2,854 & 0.04 \\
Pre-gap contiguous & 5m & 2,842 & 0.35 \\
Pre-gap contiguous & 1h & 2,063 & 1.36 \\
Post-gap contiguous & 30s & 220,456 & 8.58 \\
Post-gap contiguous & 5m & 220,244 & 9.54 \\
Post-gap contiguous & 1h & 214,610 & 9.94 \\
Exclude gap overlap (6h) & 30s & 220,456 & 8.58 \\
Exclude gap overlap (6h) & 5m & 220,246 & 9.54 \\
Exclude gap overlap (6h) & 1h & 214,616 & 9.94 \\
Exclude gap overlap (24h) & 30s & 220,456 & 8.58 \\
Exclude gap overlap (24h) & 5m & 220,246 & 9.54 \\
Exclude gap overlap (24h) & 1h & 214,616 & 9.94 \\
\bottomrule
\end{tabular}
\end{table}

\section*{S4. Secondary diagnostic layers}
\SupplementMixedEvidenceNote

\subsection*{S4.1 Timing-model misspecification}

Parametric timing fits are retained only as diagnostics. Table
\ref{tab:supp-timing-model-fit} shows that even when event probability is
approximately calibrated, the fitted conditional timing quantiles materially
misstate the early-time mass that matters for the flagship control panel.

\begin{table}[h]
\centering
\caption{Observed versus fitted timing-model diagnostics for Moltbook. Large
early-quantile residuals indicate misspecification of the seconds-to-minutes
shape under a single-family parametric timing fit.}
\label{tab:supp-timing-model-fit}
\small
\begin{tabular}{@{}lrrr@{}}
\toprule
\textbf{Statistic} & \textbf{Observed} & \textbf{Fitted} & \textbf{Residual (fit - obs)} \\
\midrule
Event probability (\%) & 9.60 & 9.91 & +0.31 \\
\(p_{10}\) (s) & 3.67 & 6.00 & +2.33 \\
\(p_{50}\) (s) & 4.55 & 39.94 & +35.39 \\
\(p_{90}\) (s) & 50.05 & 134.98 & +84.93 \\
\bottomrule
\end{tabular}
\end{table}

\subsection*{S4.2 Model-to-observable coherence}

Table \ref{tab:supp-model-observable-validation} records the in-sample
coherence check used in the current manuscript. Incidence is approximately
matched, but non-root branching is underpredicted while depth tails are
overpredicted, consistent with strong depth dependence that a coarse homogeneous
fit does not capture.

\begin{table}[h]
\centering
\caption{Model-to-observable validation: predicted versus observed incidence,
non-root branching, and depth tails.}
\label{tab:supp-model-observable-validation}
\begingroup
\setlength{\tabcolsep}{4pt}
\scriptsize
\resizebox{\linewidth}{!}{%
\begin{tabular}{@{}lrrrrrrrr@{}}
\toprule
\textbf{Group} & \textbf{Pred. inc. \%} & \textbf{Obs. inc. \%} & \textbf{Pred. branch} & \textbf{Obs. branch} & \textbf{Pred. \(\Pr(D\ge 3)\)} & \textbf{Obs. \(\Pr(D\ge 3)\)} & \textbf{Pred. \(\Pr(D\ge 5)\)} & \textbf{Obs. \(\Pr(D\ge 5)\)} \\
\midrule
Overall & 9.91 & 9.60 & 0.104 & 0.134 & 0.0109 & 0.0013 & 0.00012 & 0.00001 \\
Claimed & 20.49 & 19.23 & 0.229 & 0.274 & 0.0526 & 0.0030 & 0.00276 & 0.00000 \\
Unclaimed & 8.90 & 8.65 & 0.093 & 0.120 & 0.0087 & 0.0012 & 0.00008 & 0.00001 \\
Builder/Technical & 1.72 & 1.72 & 0.017 & 0.017 & 0.0003 & 0.0001 & 0.00000 & 0.00000 \\
Creative & 1.62 & 1.62 & 0.016 & 0.016 & 0.0003 & 0.0000 & 0.00000 & 0.00000 \\
Other & 2.27 & 2.26 & 0.023 & 0.025 & 0.0005 & 0.0001 & 0.00000 & 0.00000 \\
Philosophy/Meta & 1.81 & 1.80 & 0.018 & 0.018 & 0.0003 & 0.0002 & 0.00000 & 0.00000 \\
Social/Casual & 11.36 & 10.95 & 0.121 & 0.154 & 0.0145 & 0.0015 & 0.00021 & 0.00001 \\
Spam/Low-Signal & 0.88 & 0.88 & 0.009 & 0.009 & 0.0001 & 0.0000 & 0.00000 & 0.00000 \\
\bottomrule
\end{tabular}
}
\endgroup
\end{table}

\subsection*{S4.3 Within-thread dependence sensitivity}

Table \ref{tab:supp-dependence-robustness} reports the one-parent-per-thread
sensitivity used to bound within-thread dependence. Incidence levels move more
than conditional timing, which is consistent with the interpretation that the
main bottleneck is whether a reply happens at all rather than how slow replies
are conditional on occurring.

\begin{table}[h]
\centering
\caption{One-parent-per-thread robustness against within-thread dependence.}
\label{tab:supp-dependence-robustness}
\small
\begin{tabular}{@{}lrrrr@{}}
\toprule
\textbf{Metric} & \textbf{Primary} & \textbf{One-parent/thread} & \textbf{Abs.\ diff.} & \textbf{Rel.\ diff. \%} \\
\midrule
Reply incidence \(\Pr(\delta=1)\) & 0.09596 & 0.07213 & 0.02383 & -24.84 \\
Conditional \(t_{50}\) (s) & 4.55 & 4.51 & 0.04 & -0.93 \\
Conditional \(t_{90}\) (s) & 50.05 & 41.08 & 8.97 & -17.92 \\
Kernel half-life (diagnostic; min) & 0.68451 & 0.43601 & 0.24850 & -36.30 \\
\bottomrule
\end{tabular}
\end{table}

\subsection*{S4.4 Periodicity details}

The longest contiguous segment contains 220,461 events over 63.5175 hours.
Modulo-4-hour event-time testing on that segment gives resultant length
\(r=0.0308\), Rayleigh \(Z=209.57\), Monte Carlo
\(p=5\times 10^{-6}\), and mean phase 153.2 minutes. This rejects exact
uniformity statistically, but the concentration is extremely small. In the same
evidence base, the four-hour target in the binned-count power spectral density
diagnostics yields first-order autoregressive null-calibrated \(p\)-values
0.508, 0.501, and 0.556 for 5-, 15-, and 30-minute bins, respectively.

\begin{table}[h]
\centering
\caption{Detectability mechanics for the modulo-4-hour periodicity diagnostic.}
\label{tab:supp-periodicity-detectability}
\small
\begin{tabular}{@{}ll@{}}
\toprule
\textbf{Component} & \textbf{Value} \\
\midrule
Target period \(\tau\) & 4.0 hours \\
Segment event count & 220,461 \\
Segment duration & 63.5175 hours \\
Null Monte Carlo reps (Rayleigh) & 200,000 \\
Power simulation method & \texttt{noncentral\_chi\_square\_monte\_carlo} \\
\(\kappa\) grid & 0.0, 0.2, \ldots, 3.0 \\
Power reps per \(\kappa\) & 50,000 \\
Seed & 20260208 \\
Estimated null size at \(\kappa=0\) & 0.04964 \\
First tested \(\kappa\) with estimated power \(\ge 80\%\) & 0.2 \\
Estimated power at \(\kappa=0.2\) & 1.0 \\
\bottomrule
\end{tabular}
\end{table}

Because the \(\kappa\) grid is coarse, the resulting \(\kappa^\star=0.2\) is a
grid-crossing reference rather than a sharp detectability threshold. The
corresponding mean resultant-length mapping is \(\rho\approx 0.0995\), which is
still well above the observed \(r=0.0308\).

\subsection*{S4.5 Topic keyword sensitivity}

Deterministic keyword remapping is used only to verify the direction of the
secondary topic contrast, not to estimate semantic effects. Across the baseline
and expanded trigger dictionaries, with and without exclusions for \emph{Other}
or small categories, the Social/Casual stratum remains higher-incidence and
faster-tailed than Philosophy/Meta. Because this layer is secondary and purely
diagnostic, the detailed variant-by-variant outputs are maintained in the
reproducibility pipeline rather than repeated as a load-bearing flagship table.
